\chapter*{Résumé}
\addcontentsline{toc}{chapter}{Résumé}

Ce projet de fin d'études présente MARTA (\emph{Multi-Agent Regime-Aware Trading Architecture}), une architecture de trading algorithmique en couches construite autour d'un signal de régime de marché explicite. Le système comprend : (1) une ingénierie de variables techniques et statistiques calculées à résolution horaire sur cinq crypto-actifs (BTC, ETH, BNB, LTC, XRP) ; (2) un détecteur de régime fondé sur un modèle de Markov caché gaussien ajusté directement sur ces variables, identifiant des états de volatilité statistiquement distincts et reproductibles ; (3) une couche de délibération multi-agents dans laquelle trois agents fondés sur un grand modèle de langage (macro, technique, sentiment) forment une opinion directionnelle conditionnée par le régime, confrontée à un agent contrariant et bornée par un agent de risque GARCH sans LLM, puis agrégée de manière déterministe par un agent gestionnaire ; (4) un agent d'exécution entraîné par apprentissage par renforcement (PPO), développé et évalué comme un composant distinct qui ne consomme pas la sortie du débat ; et (5) un tableau de bord de supervision en temps réel. Un encodeur contrastif appris (TFT + InfoNCE) a été conçu initialement, puis abandonné après qu'un diagnostic a montré qu'il ne surpassait pas les variables brutes.

L'hypothèse centrale --- le conditionnement explicite par le régime améliore les résultats de trading par rapport à des approches agnostiques au régime --- a été testée avec un contrôle de bruit explicite sur deux mécanismes indépendants : des règles déterministes conditionnées par le régime et une politique PPO apprise. Dans aucun des deux cas le régime réel n'a surpassé le contrôle de bruit de manière statistiquement significative. Pour la couche multi-agents, le contexte de régime modifie effectivement les sorties des agents, et un backtest autonome sur les cinq actifs montre un Sharpe significativement supérieur à l'achat-conservation sur deux actifs sur cinq, obtenu par une exposition faible et fortement plafonnée plutôt que par un avantage directionnel ; aucune contrepartie sans régime n'étant disponible, la contribution propre du régime n'y est pas isolée. Les régimes découverts sont néanmoins des structures de volatilité réelles et reproductibles. Ce résultat négatif, obtenu avec une discipline d'évaluation rigoureuse (tests de significativité, intervalles de confiance par bootstrap par blocs, réplication hors échantillon), constitue en lui-même une contribution méthodologique du projet.

\textbf{Mots-clés :} trading algorithmique, détection de régime de marché, modèle de Markov caché, agents LLM, apprentissage par renforcement, explicabilité, validation statistique.

\chapter*{Abstract}
\addcontentsline{toc}{chapter}{Abstract}

This final-year project presents MARTA (Multi-Agent Regime-Aware Trading Architecture), a layered algorithmic trading system built around an explicit market-regime signal. The system comprises: (1) engineered technical and statistical features computed at hourly resolution on five crypto assets (BTC, ETH, BNB, LTC, XRP); (2) a Gaussian Hidden Markov Model regime detector fitted directly on those features, identifying statistically distinct, reproducible volatility states; (3) a multi-agent deliberation layer in which three large-language-model agents (macro, technical, sentiment) form regime-conditioned directional views, challenged by a contrarian agent and bounded by a non-LLM GARCH risk agent, then aggregated deterministically by a fund-manager agent; (4) a PPO reinforcement-learning execution agent, developed and evaluated as a separate component that does not consume the debate output; and (5) a real-time monitoring dashboard. A learned contrastive encoder (TFT + InfoNCE) was designed first and retired after a diagnostic showed it did not outperform the raw features.

The central hypothesis --- that explicit regime conditioning improves trading outcomes relative to regime-agnostic baselines --- was tested against an explicit noise control in two independent mechanisms: deterministic regime-conditioned rules and a learned PPO policy. In neither case did the real regime signal outperform its noise control at statistical significance. For the multi-agent layer, regime context measurably changes agent output, and a standalone five-asset backtest shows a significantly higher Sharpe ratio than buy-and-hold on two of five assets, obtained through small, tightly capped exposure rather than a directional edge; because no regime-free counterpart exists for that test, the regime's own contribution there is not isolated. The discovered regimes are nonetheless genuine, reproducible volatility structures. This negative result, reached with significance testing, block-bootstrap confidence intervals, and out-of-sample replication, is itself reported as a methodological contribution of the project.

\textbf{Keywords:} algorithmic trading, market regime detection, hidden Markov model, LLM agents, reinforcement learning, explainability, statistical validation.
